% The A2 design note.
%
% Every table in this document is \input from report/tables/, generated by
% `python -m pipeline.report` out of artifacts/tradeoffs.jsonl. No number is
% typed into this file, which is the only way the note and the repository can
% be guaranteed to agree: a figure with no ledger row behind it has nowhere to
% come from.
%
% Build:  cd report && pdflatex a2-design-note.tex && pdflatex a2-design-note.tex
% (twice, for the table of contents and the cross-references)
%
% Before building, run:
%     python -m pipeline.report
% A table whose underlying run has not happened renders as "no rows recorded",
% not as an empty tabular. That is deliberate: a note of an unfinished project
% should look unfinished rather than complete and wrong.

\documentclass[11pt,a4paper]{article}

\usepackage[margin=2.4cm]{geometry}
\usepackage{booktabs}
\usepackage{amssymb}
\usepackage{graphicx}
\usepackage{hyperref}
\usepackage[T1]{fontenc}

\newcommand{\dataset}[1]{\texttt{#1}}
% A table that may not have been generated yet. Without this, a missing file
% is a compile error in the middle of a build somebody is running an hour
% before a deadline.
\newcommand{\inputtable}[1]{%
  \IfFileExists{tables/#1.tex}{\input{tables/#1.tex}}%
  {\emph{\texttt{tables/#1.tex} not generated --- run \texttt{python -m pipeline.report}}}%
}

\title{Two-stage news recommendation on MIND and EB-NeRD:\\
       what was built, what it costs, and where it breaks}
\author{}
\date{}

\begin{document}
\maketitle

\begin{abstract}
A candidate re-ranker over a reproduced NRMS-DocVec baseline, on two news
datasets in two languages, with every design choice recorded as a comparison
rather than asserted. The system is two stages: semantic retrieval over
corrected document vectors, then a gradient-boosted re-ranker over a
forty-column feature frame that includes the baseline's own score. Every
option tried is a row in a trade-off ledger carrying both what it scores and
what it costs, and every table below is generated from that ledger.
\end{abstract}

\tableofcontents

% ===========================================================================
\section{What we built}

Two datasets, one pipeline. Every difference between MIND and EB-NeRD lives in
a registry entry --- a column map, a feature spec, a model spec --- and no
stage branches on a dataset's name. A dataset that lacks a column carries it
\emph{null} rather than not at all, so the frame has one schema and an ablation
arm is a column projection of it.

The two stages:

\begin{description}
  \item[Stage one] BM25, a FAISS flat index over corrected document vectors,
    and a linear fusion of the two. Assignment 1's work, carried forward
    unchanged and used here as candidate generation and as three feature
    columns.
  \item[Stage two] NRMS-DocVec, reproduced as a fourth retriever behind the
    same interface, and a LightGBM re-ranker over the feature frame plus the
    NRMS score --- a fifth. Both are entries in the harness's retriever table,
    so every metric, slice and bootstrap interval in this note was produced by
    one evaluation path.
\end{description}

The stacking split is chronological: NRMS fits on the earlier half of
\texttt{train}, the re-ranker on the later half, cut at a timestamp so that two
impressions of the same instant cannot land on opposite sides.

\paragraph{Chosen on \texttt{tune}, reported on \texttt{validation},
\texttt{test} scored once.}
Nothing in this note was selected on the split it is reported on.
\texttt{test} was scored in a single batch on a recorded date at a recorded
commit, by \texttt{python -m pipeline.final}, which refuses to run twice.

% ===========================================================================
\section{Options tried, and which one the registry holds}

One table per stage. Every row the sweep produced is here, losers included:
the claim is that each choice was made against alternatives, and a table of
winners is a claim without evidence. The \(\checkmark\) marks the variant the
registry actually holds --- read off the spec, not off the best row, so that a
chosen option which lost its sweep would show as losing.

\subsection{MIND}

\paragraph{Stage one retrieval.}
\inputtable{mind-retrieve}

\paragraph{The feature frame.}
\inputtable{mind-features}

\paragraph{NRMS-DocVec.}
\inputtable{mind-nrms}

\paragraph{The re-ranker.}
\inputtable{mind-rerank}

\subsection{EB-NeRD}

\paragraph{Stage one retrieval.}
\inputtable{ebnerd-retrieve}

\paragraph{The feature frame.}
\inputtable{ebnerd-features}

\paragraph{NRMS-DocVec.}
\inputtable{ebnerd-nrms}

\paragraph{The re-ranker.}
\inputtable{ebnerd-rerank}

% ===========================================================================
\section{Baseline against improved}

The headline is the paired difference \texttt{rerank} $-$ \texttt{nrms} on
\texttt{validation}, bootstrapped by impression. Paired, because the two
models rank the same impressions and an unpaired interval is dominated by
variance between impressions rather than between models --- it overlaps at any
effect size this project can produce and is a floor, not a test.

\subsection{The ablation}

Each arm is the same model with one thing removed, retrained and rescored, and
reported as a paired difference from the full model.

\paragraph{MIND.}
\inputtable{mind-ablation}

\paragraph{EB-NeRD.}
\inputtable{ebnerd-ablation}

% ===========================================================================
\section{Q9: what a causal counter costs}

Every counter the frame reads is read \emph{strictly before} the impression's
own timestamp. The leaky arm is the same model with one flag flipped: the same
store, read at a later moment.

Two rows rather than one, because the leak is largest at the cumulative window
--- where the counter runs to the end of the log --- and a single row would
quote the flattering one. Note also that the leaky \emph{sliding} window is not
a superset of the causal one: the causal hour is $[t-1h, t)$ and the leaky one
is $[t, t+1h)$, so it can read smaller, and nothing is wrong when it does.

The same question at submission time has a different answer, and the note
states it as such: a live server's click counters are the training log plus
whatever the users' own histories reveal, because the leaderboard is holding
the test period's outcomes back. Exposures it does have, strictly before each
impression.

% ===========================================================================
\section{Serving: one request, four stages}

Single-user requests through the real two-stage path, timed stage by stage
inside each request. The batched numbers do not predict these: the cost moves
from the stage that vectorises well to the stage that does not.

\paragraph{MIND.}
\inputtable{mind-serving}

\paragraph{EB-NeRD.}
\inputtable{ebnerd-serving}

\subsection{The $K$ curve}

The project's cleanest functional-against-engineering trade-off: what
retrieval depth buys in AUC and costs in p99. The AUC column comes from the
ablation's literal-cut arms, scored on \texttt{validation}; the latency column
from the serving benchmark. Same $K$, two runs, one row.

\paragraph{MIND.}
\inputtable{mind-curve-serve-k}

\paragraph{EB-NeRD.}
\inputtable{ebnerd-curve-serve-k}

\subsection{Cost per thousand queries}

Derived from measured single-core throughput at a stated SLA and a stated
price per core-hour, with the division shown. A variant whose p99 exceeds the
budget is not given a price: it does not meet the budget the price was quoted
at.

% ===========================================================================
\section{Where $10\times$ breaks}

Three axes, because they grow different stores:

\begin{description}
  \item[$10\times$ users] grows the per-user NRMS cache and the history table.
    The index does not move --- more readers do not make more articles.
  \item[$10\times$ catalogue] grows the index, the article vectors and the
    counter store: everything keyed by article.
  \item[$10\times$ QPS] grows neither. It grows cores, and it is the only one
    of the three that a bigger node does not fix.
\end{description}

The numbers, and which of the three exceeds the node first, are in
\texttt{artifacts/bench-serve-<dataset>.md}, generated by the same run as the
serving table above.

\subsection{Chunk size, from the submission runs}

\paragraph{MIND.}
\inputtable{mind-curve-chunk}

% ===========================================================================
\section{Limitations}

Stated rather than discovered by a reader:

\begin{itemize}
  \item \textbf{No position bias correction} on either dataset. Neither ships
    the position a candidate was displayed at, so it cannot be modelled and is
    not claimed to be.
  \item \textbf{MIND has no publish time.} Freshness on MIND is the log's
    first sighting of an article, which is a different quantity from age and
    is named as such in the frame.
  \item \textbf{Sessions only where shipped.} EB-NeRD has a session id and
    MIND does not; MIND's session columns are null throughout, and the
    \texttt{-session/dwell} ablation arm therefore measures nothing on MIND.
  \item \textbf{The literal-cut arm measures stage-one recall}, not the
    re-ranker. A candidate the semantic retriever would never have surfaced is
    ranked last; what that costs in AUC is a property of the retriever's
    recall@K, which is reported beside it.
  \item \textbf{The leaderboard is a different week} from \texttt{test} on both
    datasets, so a gap between the two is expected and its sign is reported
    rather than explained away.
\end{itemize}

% ===========================================================================
\section{Reproducing this}

\begin{verbatim}
  python build.py                        # every stage through rerank + evaluate
  python -m pipeline.ablation --split validation
  python -m pipeline.serve
  python -m pipeline.final               # test, once
  python -m pipeline.report              # regenerate every table above
\end{verbatim}

\noindent
See \texttt{README.md} for the credential the MIND download needs and
\texttt{ada/README.md} for the cluster placement. Every design decision that
was not obvious is recorded in \texttt{DECISIONS.md}, with what it was decided
against.

\end{document}
